\chapter{Теория}

\section{Использванные инструменты}
В ходе работы использовались такие библиотеки как NumPy,Pandas,matplotlib. \\

\subsection{NumPy}
NumPy (Numerical Python) — библиотека для быстрых вычислений с массивами чисел. Её ключевой объект — ndarray (n-мерный массив).
Главная фишка: операции выполняются над массивами целиком, без Python-циклов. Это очень ускоряет код. Помимо этого в библиотеке есть
много математических функций.

\subsection{Pandas}
Pandas — библиотека для работы с табличными и временными данными. Основные типы данных: series — столбец данных, DataFrame — таблица.
Позволяет удобно считывать .cvs файлы и анализировать их (подсчет статистических метрик, группировка записей по признаку и т.д).

\subsection{matplotlib}
Matplotlib - библиотека для построения графиков разной сложности (2D и 3D). Построение графиков происходит через pyplot — быстрый и простой интерфейс.

\endinput